\bookchapter{Chapitre 6 : Protéger votre vie numérique — sécurité et vie privée}

\chapterintro{Sécurité du smartphone = sécurité de la vie privée + sécurité du patrimoine. À l'ère où les données sont un actif, protéger votre smartphone, c'est vous protéger.}

\booksection{Écran de verrouillage : la première ligne de défense physique}

L'écran de verrouillage est la première ligne de défense de votre smartphone, et aussi la plus facilement négligée.

\booksubsection{Classement des niveaux de sécurité :}

\begin{enumerate}
  \item \textbf{Mot de passe/code PIN}: Sécurité la plus élevée, garantie ultime (quand vous oubliez l'empreinte/le visage).
  \item \textbf{Reconnaissance d'empreinte digitale}: Équilibre parfait entre sécurité et commodité.
  \item \textbf{Reconnaissance faciale}: Le plus pratique, mais la sécurité dépend de la technologie (3D Face ID sécurisé, 2D facilement trompé par une photo).
  \item \textbf{Déverrouillage par schéma}: Pratique mais laisse facilement des traces devinables.
\end{enumerate}

\booksubsection{Configurer un verrouillage fort :}

\begin{itemize}
  \item Configurez impérativement un mot de passe complexe (n'utilisez pas votre date de naissance, 123456).
  \item Activez la biométrie.
  \item Réglez le délai de verrouillage sur « Immédiat » ou le plus court possible.
\end{itemize}

\booksection{« Localiser mon smartphone » : l'outil antivol ultime}

\textbf{Il faut l'activer !}Il permet de :

\begin{itemize}
  \item \textbf{Localiser}: Où est le smartphone perdu, le retrouver sur la carte.
  \item \textbf{Jouer un son}: Faire sonner le smartphone, pratique pour le chercher à la maison.
  \item \textbf{Mode perdu}: Verrouiller à distance le smartphone et afficher un numéro de contact.
  \item \textbf{Effacement à distance}: Dernier recours, effacer toutes les données (mais après effacement, la localisation n'est plus possible).
\end{itemize}

Allez dans [Paramètres], Apple cherchez « Localiser mon iPhone », Android cherchez « Localiser mon appareil » ou un service similaire de la marque, assurez-vous qu'il est activé (vert).

\booksection{Mise à jour système : appliquer les correctifs, boucher les failles}

Le but principal des mises à jour système est de corriger les failles de sécurité, comme reboucher les trous d'un mur. Donc, quand vous voyez une notification de mise à jour système, essayez de mettre à jour le plus vite possible.

Bien sûr, les mises à jour de version majeure (comme iOS 17 vers iOS 18) peuvent avoir des problèmes de compatibilité, vous pouvez attendre quelques jours avant de mettre à jour.

\booksection{Identifier et prévenir les fraudes téléphoniques}

Avec le développement rapide des technologies IA, le modèle cognitif traditionnel du « voir pour croire » est totalement caduc. Les escrocs utilisent désormais des technologies IA comme le Deepfake, passant de l'arnaque « à l'aveugle » à la criminalité « ciblée de précision ».

\booksubsection{Les types d'arnaques les plus courants à l'ère de l'IA}

\textbf{(1) Arnaque par usurpation de visage et de voix par IA}\par

Les malfaiteurs collectent illégalement des photos haute définition publiées par la cible sur les réseaux sociaux ou quelques secondes d'échantillon vocal, puis utilisent la technologie de deepfake IA pour cloner une « vidéo de proche » ou une « voix » suffisamment convaincante. Scénarios courants :\textbf{accident de voiture urgent, prêt d'urgence, enlèvement}, etc. Ce type d'arnaque exploite les liens affectifs familiaux et amicaux, avec une\textbf{force de tromperie extrêmement élevée}.

\textbf{(2) Usurpation d'autorité / arnaque au service client}\par

La technologie IA peut simuler avec précision la voix et le ton du personnel du service client officiel, des autorités judiciaires et policières, et même synthétiser des captures d'écran de virement très réalistes, des documents administratifs électroniques, des mandats d'arrêt, etc., en\textbf{créant un climat de panique ou une pression d'autorité}, incitant la victime à transférer de l'argent ou à activer le partage d'écran, volant ainsi les informations du compte.

\textbf{(3) Arnaque IA à l'investissement financier}\par

Via des algorithmes IA, falsification de captures d'écran de rendements « garantis sans risque », d'interfaces de logiciels de trading quantitatif virtuels ou de rapports d'analyse d'investissement de « niveau maître », avec la promesse de rendements élevés comme appât, vente de services d'investissement onéreux ou incitation des victimes à transférer des fonds vers des comptes de trading sous leur contrôle, pour\textbf{finalement disparaître avec l'argent}.

\booksubsection{Comment les démasquer — l'approche du « miroir révélateur »}

L'IA peut cloner la voix et le visage, mais elle n'est pas invulnérable, il y a plusieurs failles exploitables :

\begin{itemize}
  \item \textbf{Méthode de vérification par action}: Lors d'un appel vidéo, demandez activement à l'interlocuteur d'exécuter des actions consécutives prédéfinies, par exemple « passez rapidement la main devant votre visage », « montrez le contour de votre oreille de profil » ou « clignez rapidement 5 fois de suite ». Actuellement, la technologie de deepfake IA a du mal avec les occlusions dynamiques complexes, les scènes de mouvement rapide et les changements d'ombre,\textbf{avec des artefacts très probables comme des saccades, des distorsions du visage ou des bords défectueux}.
  \item \textbf{Méthode de vérification par mot de passe secret}: Établissez et utilisez des questions de vérification privées prédéfinies, comme « Quel plat avez-vous commandé lors de notre dernier dîner ensemble », « Quel a été le premier cadeau que je vous ai offert » et autres détails connus seulement des deux parties. Si l'interlocuteur hésite visiblement ou répond à côté,\textbf{mettez immédiatement fin à l'appel et activez le protocole de sécurité}.
  \item \textbf{Méthode de double vérification}: Respectez strictement le principe des «\textbf{quatre non et une vérification}» publié par le Bureau de la cybersécurité du ministère de la Sécurité publique — ne pas croire aveuglément les messages inconnus, ne pas transférer d'argent à la légère, ne pas divulguer d'informations personnelles sensibles, ne pas diffuser de contenu suspect. Pour toute transaction financière, vérifiez impérativement en\textbf{rappelant le numéro d'origine connu ou en confirmant en personne}. Pour les virements, il est recommandé de choisir le mode « crédit sous 24 heures », laissant une fenêtre de temps suffisante pour identifier une arnaque potentielle et récupérer les fonds.
\end{itemize}

\booksubsection{Schémas traditionnels courants (restez vigilants) :}

\begin{itemize}
  \item \textbf{Loterie/cashback par fausse tâche}: Appât, d'abord un petit gain puis une grosse arnaque.
  \item \textbf{Arnaque sentimentale}: « Pig butchering scam », cultiver une relation puis escroquer.
  \item \textbf{Usurpation d'identité d'un proche}: Vol de compte puis demande d'argent (encore plus difficile à distinguer à l'ère de l'IA).
\end{itemize}

\booksubsection{Protection intégrée du smartphone :}

\begin{itemize}
  \item Blocage des appels indésirables.
  \item Protection contre le vol de codes de vérification.
  \item Détection de sécurité des liens SMS.
\end{itemize}

\booksubsection{Formule anti-arnaque :}

\textbf{Ne pas croire aveuglément, ne pas divulguer, ne pas cliquer, ne pas transférer, en cas de doute, vérifiez d'abord !}Face à une « urgence », prenez 5 minutes pour vous calmer, presque toutes les arnaques se démasquent à ce moment-là.

\booksection{Résister au « doxxing » et aux attaques d'ingénierie sociale : couper la chaîne de fuite de données privées}

Le principe central du « doxxing » (attaque d'ingénierie sociale) est de\textbf{collecter systématiquement, analyser et assembler les fragments numériques personnels pour construire un profil complet de la personne}. La clé de la prévention est de tracer des lignes rouges claires pour la vie privée personnelle et d'établir un système de protection multicouche.

\booksubsection{Anonymisation des photos}

Les méthodes de base de nettoyage EXIF se trouvent au chapitre 2 « Gestion des images et du système de fichiers ». Avant de partager des photos sur les réseaux sociaux, il faut aller plus loin\textbf{supprimer complètement les métadonnées de la photo (informations EXIF)}, y compris la date de prise de vue, le modèle de l'appareil, les coordonnées GPS et autres données sensibles. Les informations clés contenues dans la photo comme le\textbf{numéro de porte, la plaque d'immatriculation, le code-barres du billet, le numéro de carte d'identité}doivent être strictement pixellisées ou floutées, évitant de révéler involontairement vos déplacements.

\begin{itemize}
  \item iOS : Dans [Paramètres] > [Confidentialité et sécurité] > [Appareil photo], vous pouvez désactiver « Localisation précise » ; lors du partage, choisir « Uniquement la photo actuelle » permet de supprimer les métadonnées.
  \item Android : Certaines applications de galerie (comme Google Photos) vous demandent si vous souhaitez conserver les informations de localisation lors du partage.
  \item Outils universels : ExifTool, Metapho, etc. permettent de voir/éditer les EXIF.
\end{itemize}

\booksubsection{Restriction des autorisations d'application}

Les méthodes de gestion des permissions de base sont détaillées au chapitre 4 « Gestion des permissions ». Dans un scénario de « protection contre le doxxing », il faut en plus vérifier et fermer régulièrement les permissions des applications de e-commerce et d'outils\textbf{micro, les contacts, la localisation}, empêchant les applications d'« écouter » vos conversations en arrière-plan ou de collecter illégalement votre réseau relationnel ; les applications inactives depuis longtemps mais occupant encore des autorisations élevées doivent être désinstallées rapidement, réduisant les points de fuite de données privées à la source.

\begin{itemize}
  \item Quand vous voyez une demande d'« autorisation des contacts », posez-vous la question : cette application a-t-elle vraiment besoin de connaître tous mes contacts ?
  \item Une application non utilisée depuis trois mois, désinstallez-la sans hésiter. Les autorisations se resserrent avec le temps, pas l'inverse.
\end{itemize}

\booksubsection{Sécurisation des appareils connectés}

Les caméras connectées, serrures connectées, enceintes connectées et autres appareils IoT doivent impérativement avoir leur\textbf{mot de passe de connexion par défaut modifié immédiatement}lors de la première utilisation, en configurant un mot de passe indépendant de haute intensité comprenant majuscules, minuscules, chiffres et caractères spéciaux ; les fonctions d'accès à distance peu utilisées doivent rester désactivées, vérifiez régulièrement les mises à jour du firmware et corrigez les failles de sécurité en temps utile,\textbf{pour empêcher les pirates d'infiltrer et d'espionner l'intimité du foyer}.

\begin{tcbwarn}{Avertissement basé sur un cas réel}
En 2022, la police d'une localité a signalé que des pirates avaient cracké les mots de passe par défaut de caméras de surveillance domestiques d'une marque (beaucoup d'utilisateurs ne les avaient jamais modifiés), volant les images de chambres de milliers de foyers et les revendant en ligne.\textbf{Ne pas changer le mot de passe par défaut = laisser la porte de chez soi grande ouverte}.

\end{tcbwarn}

\booksection{Protection de la vie privée : vos données sont votre patrimoine}

Les photos, historiques de discussion et traces de localisation dans votre smartphone sont en essence vos « actifs numériques ». Une fois divulgués, au mieux du spam publicitaire, au pire du doxxing, de la fraude voire du chantage.

\booksubsection{Principe de « minimisation » des autorisations d'application :}

\begin{itemize}
  \item \textbf{Localisation}: VTC/carte en « Toujours autoriser » ; les autres applications en « Autoriser uniquement pendant l'utilisation ».
  \item \textbf{Carnet de contacts}: Sauf nécessité (comme ajouter un ami WeChat), interdire systématiquement.
  \item \textbf{Album / Toutes les photos}: Depuis iOS 14+, vous pouvez autoriser « Photos sélectionnées » ; Android 13+ prend également en charge l'autorisation fine.
  \item \textbf{Micro / Caméra}: Désactivés par défaut, activez au besoin. En cas de suspicion d'enregistrement à votre insu, faites glisser le centre de contrôle pour vérifier la présence d'un « petit point vert/petit point orange ».
  \item \textbf{Suivi en arrière-plan}: Depuis iOS 14.5+, les applications doivent obligatoirement demander l'autorisation avant d'utiliser l'IDFA ; Android permet de désactiver l'« ID publicitaire ».
\end{itemize}

\booksubsection{Gestion pratique de la vie privée :}

\begin{itemize}
  \item \textbf{Rapport de confidentialité (iOS)}: [Paramètres] > [Confidentialité et sécurité] > [Rapport de confidentialité des apps], voir quelles applications accèdent subrepticement à votre localisation, appareil photo, micro.
  \item \textbf{Historique d'utilisation des autorisations (Android)}: [Paramètres] > [Vie privée] > [Historique d'utilisation des autorisations], voir qui a accédé à des autorisations sensibles et quand.
  \item \textbf{Limiter le suivi publicitaire}: Désactivez « ID publicitaire » ou « Limiter le suivi publicitaire » dans les paramètres système, pour réduire le « harcèlement par publicités ciblées ».
\end{itemize}

\booksection{Liste de « premiers secours » après perte / vol}

Ne paniquez pas si vous perdez votre smartphone, suivez cet ordre :

\begin{enumerate}
  \item \textbf{Geler immédiatement la carte SIM}: Contactez l'opérateur pour la suspendre, évitant qu'elle serve à recevoir des codes de vérification.
  \item \textbf{Verrouillage à distance}: Via le service cloud (iCloud / cloud de marque), activez le « Mode perdu ».
  \item \textbf{Modifier les mots de passe importants}: Changez immédiatement les mots de passe de tous les comptes importants connectés au smartphone (paiement, banque, messagerie).
  \item \textbf{Déposer plainte}: Allez au commissariat, notez le numéro IMEI. Facilite aussi la preuve d'identité si le smartphone est retrouvé.
  \item \textbf{Prévenir vos proches}: Pour empêcher les escrocs d'emprunter votre identité pour demander de l'argent.
  \item \textbf{(En dernier recours) Effacement à distance}: N'effectuez cette opération qu'en cas de perte définitive, la localisation devient inactive après effacement.
\end{enumerate}

\booksection{Sécurité de la carte SIM et de l'eSIM}

\begin{itemize}
  \item \textbf{Configurer un PIN de carte SIM}: Après activation, à chaque redémarrage ou changement de carte, le code PIN est requis, empêchant de « retirer la carte et changer de smartphone » pour contourner le verrouillage.
  \item \textbf{Avantages de l'eSIM}: Pas de carte physique à retirer, sécurité plus élevée ; prend en charge la bascule entre plusieurs numéros, très pratique en voyage.
  \item \textbf{Risque de réémission de carte}: Méfiez-vous de l'arnaque de « réémission de carte par l'opérateur » — l'escroc se fait passer pour vous pour faire réémettre la carte SIM, prenant ainsi le contrôle de votre numéro de téléphone. Configurer « la réémission nécessite la présence personnelle du titulaire avec sa carte d'identité en boutique » est la défense la plus efficace.
\end{itemize}

\begin{tcbwarn}{Ce qu'il ne faut absolument pas faire}
\begin{itemize}
  \item Utiliser des mots de passe faibles comme date de naissance / 123456 / qwerty ;
  \item Envoyer votre code de vérification à qui que ce soit (même au « service client ») ;
  \item Cliquer sur des liens inconnus dans les SMS, surtout ceux commençant par 00 (numéros internationaux) ;
  \item Se connecter à la banque en ligne ou effectuer des paiements en Wi-Fi public ;
  \item Activer le « partage d'écran » ou le « contrôle à distance » pour un inconnu.
\end{itemize}

\end{tcbwarn}

\begin{tcbpractice}{Essayez d'activer et de simuler un « verrouillage à distance »}
\begin{enumerate}
  \item Assurez-vous que la fonction « Localiser mon smartphone » est activée.
  \item Sur le navigateur d'un ordinateur ou un autre smartphone, connectez-vous au compte de service cloud correspondant (iCloud.com ou le site du service cloud de la marque).
  \item Trouvez la fonction « Localiser l'appareil », essayez d'envoyer une commande « Jouer un son » ou de simuler une commande « Verrouiller » sur votre propre smartphone, pour vous familiariser avec le processus complet.
\end{enumerate}

\end{tcbpractice}
